% Pandoc template for the fed_fraud documents.
% Based on proposal_gen/templates/detailed_proposal_template/detailed_1_classic.tex
% (Classic report): centered title page, table of contents, ruled section
% headings and the Geodesic house footer. Fonts are Latin Modern (lmodern).
% Two layouts: the full report (title page and contents) and, with
% -V brief, a compact one-page-brief header in the same style.
\documentclass[letterpaper,11pt]{article}
\usepackage[T1]{fontenc}
\usepackage[utf8]{inputenc}
\usepackage{lmodern}
\usepackage[letterpaper,top=0.9in,bottom=1.2in,left=1in,right=1in,footskip=0.6in]{geometry}
\usepackage[table]{xcolor}
\usepackage{geodesicbrand}
\usepackage{amsmath,amssymb}
\usepackage{enumitem}
\usepackage{array}
\usepackage{calc}
\usepackage{longtable}
\usepackage{makecell}
\usepackage[most]{tcolorbox}
\usepackage{titlesec}
\usepackage{needspace}
\usepackage{etoolbox}
\usepackage{hyperref}
\hypersetup{colorlinks=true,urlcolor=coolblack,linkcolor=coolblack,filecolor=geoblue,
  pdftitle={Army Corps of Engineers (Civil Works)},pdfauthor={\geofirm}}
\urlstyle{same}

\setlist{itemsep=2pt,leftmargin=1.4em}
\setlength{\parindent}{0pt}
\setlength{\parskip}{6pt}
\providecommand{\tightlist}{\setlength{\itemsep}{0pt}\setlength{\parskip}{0pt}}
\setcounter{secnumdepth}{0}

% Section style: house rule under each heading
\titleformat{\section}{\Large\bfseries\color{coolblack}}{\thesection}{0.6em}{}[{\color{coolblack}\titlerule[0.5pt]}]
\titleformat{\subsection}{\large\bfseries\color{geomid}}{\thesubsection}{0.6em}{}
\titleformat{\subsubsection}{\normalsize\bfseries\color{geomid}}{\thesubsubsection}{0.6em}{}
\titlespacing*{\section}{0pt}{10pt}{6pt}

% Tables are written by latex/filters.lua (navy header, banded rows)
\setlength{\LTpre}{6pt}\setlength{\LTpost}{6pt}

\begin{document}

% BRIEF HEADER
\thispagestyle{geodesic}
\begin{minipage}[c]{0.13\textwidth}\includegraphics[width=\linewidth]{icon.png}\end{minipage}\hfill
\begin{minipage}[c]{0.84\textwidth}
  {\textsc{\large \geofirm}}\par\vspace{2pt}
  {\small\color{geogray}Federal fraud expansion\enspace\textperiodcentered\enspace When: \compiletime}
\end{minipage}\par\vspace{6pt}
{\color{coolblack}\rule{\linewidth}{0.3mm}}\par\vspace{2pt}
{\LARGE\bfseries\color{coolblack} Army Corps of Engineers (Civil Works)}\par
\vspace{2pt}{\large\color{geomid} One-page brief}\par
\vspace{-2pt}{\color{coolblack}\rule{\linewidth}{0.3mm}}\par\vspace{4pt}
\pagestyle{geodesic}

\textbf{Priority:} 14 of 33 (score 3.2). \textbf{FY2026 outlays:} \$28B
(USAspending). \textbf{As of:} 2026-10-27. Public data only.

\hypertarget{bottom-line}{%
\section{Bottom line}\label{bottom-line}}

The Corps spends almost entirely through construction and dredging
contracts, many of them small business set-asides, and it ramps up
sharply after disasters. That makes it a near copy of the DoD contractor
screen at a smaller scale. Bid rigging among a few regional dredging and
construction firms is the second pattern.

\hypertarget{exposure}{%
\section{Exposure}\label{exposure}}

\begin{itemize}
\tightlist
\item
  \$28B in outlays and \$77B in budgetary resources in FY2026, boosted
  by infrastructure and disaster supplemental money.
\item
  Set-aside fronts, bid rigging and inflated change orders are recurring
  case types for the Army CID and DOJ's Procurement Collusion Strike
  Force.
\end{itemize}

\hypertarget{public-data-we-can-use}{%
\section{Public data we can use}\label{public-data-we-can-use}}

\begin{itemize}
\tightlist
\item
  \textbf{FPDS / USAspending:} awards, set-aside type, number of offers,
  modifications.
\item
  \textbf{SAM.gov} entities and exclusions; SBA certification lists.
\item
  \textbf{DOJ and Army CID} case releases.
\end{itemize}

\hypertarget{detection-plays}{%
\section{Detection plays}\label{detection-plays}}

\begin{enumerate}
\def\labelenumi{\arabic{enumi}.}
\tightlist
\item
  \textbf{Front companies:} the DoD set-aside screen (see \url{dod.pdf})
  applied to Corps districts.
\item
  \textbf{Collusion signals:} few bidders per award, stable win
  rotation, and shared officers among ``competitors'' in a district.
\item
  \textbf{Change-order growth:} contracts whose value grows far beyond
  award compared with peer projects.
\end{enumerate}

\hypertarget{first-prototype}{%
\section{First prototype}\label{first-prototype}}

Run the DoD contractor integrity screen on Corps awards and add a
district-level bid rotation view.

\hypertarget{limits}{%
\section{Limits}\label{limits}}

Bid amounts for losing offers are not in FPDS; collusion tests rely on
number of offers and win patterns.

\hypertarget{who-to-take-it-to}{%
\section{Who to take it to}\label{who-to-take-it-to}}

USACE Directorate of Contracting, Army CID Major Procurement Fraud Unit,
DoD OIG, and DOJ's Procurement Collusion Strike Force.

\hypertarget{sources}{%
\section{Sources}\label{sources}}

\begin{itemize}
\tightlist
\item
  \href{https://www.usaspending.gov/agency}{USAspending agency profiles
  (FY2026 outlays)}
\item
  \href{https://www.justice.gov/atr/procurement-collusion-strike-force}{DOJ
  Procurement Collusion Strike Force}
\end{itemize}

\end{document}
